% ============================================================
% ch. 15 Glossary — Arabic edition (Phase 28R, unit 16)
% Transcribed from the accepted edition text layer
% (book-ar-fixed.pdf b78c2a28…, accepted pp. 180–190, PRIMARY per
% the unit 14 method ruling) with the recovery reference
% (recovery-r0/text/15-glossary.txt) as cross-check and spacing
% arbiter, no retranslation. Upstream glossary/glossary.tex gives
% structure only (description environment, itemize/enumerate
% sub-structures, cite keys, siunitx values, labels).
%
% Structure (verified against accepted renders pp. 180/182/183/184/185/187):
%   - quoting intro, then ONE description environment, 72 \item[term]
%     entries; term = Arabic (bold via description label) + English
%     parenthetical.
%   - Baking entry: 7-item ENUMERATE (numbered 1.–7. in the accepted
%     edition, Western numerals).
%   - Float test: two unmarked \item[] paragraphs (النتيجة الإيجابية /
%     النتيجة السلبية) + closing paragraph OUTSIDE the list.
%   - P/L Value: two unmarked \item[] groups (P/L breakdown, low/high)
%     + closing paragraph OUTSIDE the lists.
%   - Upstream \emph{} italics kept inside their \textenglish{} islands
%     (alveogram, oven spring, tight crumb, open crumb, pain de mie);
%     alveoli translated in parens per accepted (الحويصلات);
%     Brühstück "boil piece" translated «القطعة المسلوقة»; "Also see
%     scald" rendered as the entry cross-reference انظر أيضًا مدخل
%     السكب المغلي (accepted wording).
%   - Islands: \qty{} rides \LR{}; pH values ride \LR{pH 6.0}; bare
%     digit runs per accepted layer (0/14/7 bare; 30/60 per-token
%     \LR{}; 4.2 bare per the unit 14 same-sentence precedent);
%     ratios \LR{1:1}; cites ~\cite{wheat+amylase+singh,
%     innograin+enzymes,reichenberger+amylase} render [1],[2],[3] vs
%     accepted [55, 33, 50] (by-design unit bibliography numbering,
%     R3 assembly regenerates true numbering).
% ============================================================

% Unnumbered chapter per the accepted edition (TOC entry «المسرد 180»
% carries no chapter number; chapter head on accepted p. 180 has no
% number, rule spans full width): \addchap, not \chapter.
\addchap{المسرد}%
\label{ch:Glossary}

\begin{quoting}
يقدّم هذا المسرد تعريفات وشرحًا للمصطلحات التي تستخدم بكثرة في صناعة الخبز. إن فهم هذه المصطلحات ضروري للخبازين المبتدئين وذوي الخبرة على حدٍّ سواء، لمن يطمح إلى إتقان فن صناعة الخبز وعلمه معًا. وقد رُتِّب هذا المسرد ترتيبًا أبجديًا لسهولة الرجوع إليه.
\end{quoting}

\begin{description}

\item[حمض الخليك \textenglish{(Acetic Acid)}]
نوع من الأحماض العضوية تُنتجه بكتيريا حمض اللبنيك المتخمِرة بشكل مغاير \textenglish{(heterofermentative)} وبكتيريا حمض الخليك أثناء التخمير. يمنح الخبز الحامض طعمه الحامض المميز، ويساعد على حفظ الخبز بخفض الرقم الهيدروجيني \textenglish{(pH)}. وطعم حمض الخليك صاحب بروفايلٍ أقرب إلى الخل.

\item[جرّة أليكوت \textenglish{(Aliquot jar)}]
قطعة صغيرة من العجين تُستخرج بعد بناء قوة العجين الأولية. تُستخدم جرّة أليكوت لمتابعة تقدّم تخمير العجين. من المهم التأكد من أن حرارة ماء العجين داخل الجرّة تطابق حرارة الغرفة لديك لكي تكون القراءات دقيقة. وانتبه إلى أن جرّة أليكوت قد لا تكون فعّالة إذا كانت تذبذبات حرارة مطبخك كبيرة؛ لأن عيّنة العجين الصغيرة داخل الجرّة ترتفع حرارتها أو تنخفض أسرع من كتلة العجين الرئيسية، ما قد يضرّ بقدرتها على متابعة التخمير بدقة. ومن الجوهري استخدام وعاء جرّة أسطواني الشكل حتى يمكنك الحكم بشكل سليم على زيادة حجم العجين.

\item[دقيق لجميع الأغراض \textenglish{(All Purpose Flour)}]
دقيق عام متوازن لصنع الخبز والكيك معًا. في ألمانيا هذا الدقيق من النوع \LR{550}.

\item[ألفا أميلاز \textenglish{(Alpha-amylase)}]
نوع من الأميلاز يفكك جزيئات النشا إلى شظايا أقصر، منتجًا المالتوز وبعض الغلوكوز.

\item[الألفيوغراف \textenglish{(Alveograph)}]
جهاز يُستخدم أساسًا في تقييم جودة الخَبز لدقيق القمح. يقيّم الألفيوغراف الخصائص الريولوجية للعجين، وقابليته للتمدد ومقاومته للمدّ تحديدًا، بنفخ قطعة من العجين كبالونة حتى تنفجر. والرسم الناتج، أو \textenglish{\emph{alveogram}}، يعرض منحنى يمثل التوازن بين مرونة العجين وقابليته للتمدد. وتوفر المعايير المشتقة من المنحنى، مثل \textenglish{P} (الضغط اللازم لنفخ العجين) و\textenglish{L} (قابلية العجين للتمدد)، إرشادات لا تقدَّر بثمن للخبازين والمطاحن بشأن الأداء المحتمل للدقيق في صناعة الخبز. فبتحليل الألفيوغرام يمكن للمهنيين اتخاذ قرارات رشيدة بشأن ملاءمة دقيقٍ ما لتطبيقات خبزٍ معينة، وكذلك الحاجة المحتملة لخلطه بدقائق أخرى.

\item[الحويصلات \textenglish{(Alveoli)}]
(المفرد: حويصلة \textenglish{alveolus}) الجيوب الصغيرة التي تكوّن لُبّ الخبز، وتتخلّق نتيجة احتجاز مصفوفة الغلوتين لسهام ثاني أكسيد الكربون داخلها.

\item[الأميلاز \textenglish{(Amylase)}]
إنزيم يفكك النشويات إلى سكريات أبسط، ميسّرًا عملية التخمير في صناعة البيرة والخبز. عند صناعة البيرة تُحفظ حرارة الشراب مددًا طويلة عند درجات حرارة معينة لضمان تفكك معظم النشا إلى سكريات، تستهلكها ثم الكائنات الدقيقة أثناء عملية التخمير.

\item[الأوتولايز \textenglish{(Autolyse)}]
عملية يُخلط فيها الدقيق مع الماء ثم يُترك ليرتاح قبل إضافة بقية المكوّنات. ينشّط هذا إنزيمات مثل الأميلاز والبروتياز. وبفعل ذلك يقصُر زمن التخمير الأول، وتكون الرغيفة النهائية ذات خصائص أفضل؛ فيتحسن اسمرار الرغيفة ويلين لُبّها. يوصى بالأوتولايز عند استخدام نسبة مرتفعة من البادئة لتلقيح العجين (فوق \LR{\qty{20}{\percent}}). ويمكن أن يكون البديل الأسهل هو الفيرمنتولايز.

\item[البكتيريا \textenglish{(Bacteria)}]
كائنات دقيقة وحيدة الخلية موجودة في أشكال وموائل متنوعة. تؤدي أدوارًا حاسمة في عمليات طبيعية متعددة، لا سيما في إعداد الغذاء مثل تخمير الخبز الحامض. وتلعب بكتيريا حمض اللبنيك وحمض الخليك خصوصًا دورًا محوريًا في عملية الخبز الحامض، مسهمةً في مذاقه وقوامه المميزين. بعض البكتيريا نافعة، تساعد في الهضم أو تنتج فيتامينات، وبعضها الآخر ضار وقد يسبب الأمراض.

\item[الحساب الخبازي \textenglish{(Baker's Math)}]
الحساب الخبازي نظام قائم على النسب لمشاركة الوصفات، يجعلها قابلة للتوسيع بسهولة. يقوم على الوزن الكلي للدقيق في الوصفة، حيث يُقسم وزن كل مكوّن على وزن الدقيق ليُعطي نسبة مئوية. فمن دقيق بوزن \LR{\qty{500}{\gram}} يمكن استخدام \LR{\qty{60}{\percent}} ماء (\LR{\qty{300}{\gram}})، و\LR{\qty{10}{\percent}} بادئة (\LR{\qty{50}{\gram}})، و\LR{\qty{2}{\percent}} ملح (\LR{\qty{10}{\gram}}).

\item[النسبة الخبازية \textenglish{(Baker's percentage)}]
انظر الحساب الخبازي.

\item[الخَبز \textenglish{(Baking)}]
الخطوة الأخيرة التحويلية في صناعة الخبز، حيث يتعرض العجين لدرجات حرارة مرتفعة تُحدث سلسلة من التفاعلات الكيميائية والمادية التي تنتهي برغيف خبز مكتمل. خلال مرحلة الخَبز:
\begin{enumerate}
    \item نشاط الخميرة والانتفاخ في الفرن: في الطور الأول من الخَبز ترتفع الحرارة داخل العجين، فيزداد نشاط الخميرة. ويؤدي ذلك إلى إنتاج سريع لثاني أكسيد الكربون، مما يحقق ما يسميه الخبازون الانتفاخ في الفرن \textenglish{(\emph{oven spring})}، أي النمو السريع للرغيف.
    \item تخثار البروتينات: مع استمرار ارتفاع الحرارة تبدأ البروتينات في العجين، الغلوتين أساسًا، بالتخثار والتصلب، وهي ما يمنح الخبز بنيته.
    \item جيلاتنة النشا: تمتص النشويات الماء وتنتفخ، ثم تجلتن في النهاية. وتسهم هذه العملية في بنية لُبّ الخبز.
    \item التكرمل وتفاعل مايار: يسمُر قشر الخبز بفعل تفاعلين رئيسيين: تكرمل السكريات، وتفاعل مايار بين الأحماض الأمينية والسكريات المُرجِعة. وهذا لا يؤثر في المظهر فحسب، بل يمنح الخبز أيضًا نكهته ورائحته المميزتين.
    \item تبخر الأحماض: بعض الأحماض الناتجة أثناء التخمير يتبخر عند درجات حرارة معينة أثناء الخَبز. وهذا التبخر قد يؤثر في بروفايل النكهة النهائية للخبز، فيجعله أقل حامضية من العجين غير المخبوز. وبمدّ زمن الخَبز تنخفض تركُّز الأحماض وقد يفقد العجين بعض حامضيته.
    \item تبخر الرطوبة: يتحول الماء في العجين إلى بخار ويبدأ بالتبخر. ويسهم البخار في الانتفاخ في الفرن ويساعد أيضًا على جيلاتنة النشا.
    \item تكوّن القشرة: تجف الطبقة الخارجية من العجين وتتصلب لتشكل قشرة تعمل حاجزًا واقيًا يحفظ رطوبة اللُّبّ الداخلي.
\end{enumerate}

\item[سلة التخمير \textenglish{(Banneton)}]
سلة من الخوص تُستخدم لتشكيل العجين ودعمه خلال تخميره النهائي. تُصنع سلال التخمير عادة من الخيزران أو لبّ الخشب. والحل اليدوي البديل هو استخدام وعاء توضع داخله منشفة مطبخ. وأثناء راحة العجين داخل سلة التخمير يجف سطح العجين ويصبح أسهل في التشريط قبل الخَبز.

\item[الباسيناج \textenglish{(Bassinage method)}]
تقنية من تقنيات صناعة الخبز تعتمد على الإضافة المتدرجة للماء إلى العجين. يُخلط العجين أولًا عند مستوى ترطيب أدنى، مما يسمح لروابط الغلوتين بالتشكل بفعالية أكبر. وبعد تشكل هذه البنى الغلوتينية يُدمَج الماء الإضافي تدريجيًا عبر عجنٍ إضافي. تعزز هذه الطريقة قابلية العجين للتمدد، وهي مفيدة خصوصًا عند العمل بدقائق منخفضة الغلوتين. فباستخدام طريقة الباسيناج يستطيع الخباز الحصول على عجين قوي وقابل للتمدد معًا.

\item[الراحة على السطح \textenglish{(Bench Rest)}]
فترة راحة قصيرة تُمنح للعجين بعد التشكيل الأولي، تتيح للغلوتين أن يرتخي قليلًا وتجعل التشكيل أسهل. يترك أغلب الناس العجين على السطح من \LR{10} دقائق حتى ساعة. وتصبح هذه الراحة مهمة خصوصًا عند صنع عجائن البيتزا. فبدون راحة مطوّلة على السطح يصير العجين مرونًا أكثر من اللازم ولا يمكن تشكيله.

\item[بيتا أميلاز \textenglish{(Beta-amylase)}]
إنزيم يفكك أكثر شظايا النشا الناتجة عن الألفا أميلاز إلى المالتوز.

\item[دقيق الخبز \textenglish{(Bread Flour)}]
دقيق مثالي لصناعة الخبز الحامض. يحتوي نسبة أعلى من الغلوتين، ولهذا يستطيع أن يختمر مدة أطول.

\item[القطعة المغلية \textenglish{(Brühstück)}]
تقنية خبز ألمانية تشبه السكب المغلي. وتعني حرفيًا «القطعة المسلوقة». يُسكب الماء الساخن أو المغلي فوق دقيق كامل الحبة أو حبوب مطحونة، ثم يُترك ليبرد ويُخلط مع العجين الرئيسي. تساعد هذه العملية على احتفاظ العجين بالرطوبة، وقد تعزز نكهة الخبز النهائي وقوامه. انظر أيضًا مدخل السكب المغلي.

\item[التخمير الأول \textenglish{(Bulk Fermentation)}]
فترة النمو الأولى بعد خلط كل المكوّنات. يُترك العجين عادة لينمو حتى يزداد إلى حجم معين. وتعتمد نسبة الزيادة في الحجم على الدقيق المستخدم. فعند الخَبز بدقيق القمح فإن كمية الغلوتين في الدقيق هي العامل الحاسم: كلما زاد الغلوتين في دقيقك (البروتين) طال مدة التخمير الأول. والتخمير الأول الأطول يحسّن نكهة الخبز النهائي وقوامه، فيصبح أكثر حامضية ونعومة. يمكنك أن تستهدف زيادة في حجم عجينك بمقدار \LR{\qty{25}{\percent}} ثم تزيد هذه النسبة تدريجيًا لتجد النقطة المثالية لدقيقك؛ فهذا الأمر يعتمد على الدقيق إلى الدقيق بدرجة عالية. وعند استخدام دقيق منخفض الغلوتين مثل الجاودار عليك الحذر، لأن التخمير الأطول قد يخلق عجينًا لزجًا أكثر من اللازم ينهار ولم يعد يحفظ شكله.

\item[دقيق الكعك \textenglish{(Cake Flour)}]
دقيق الكعك دقيق فاتح ناعم الطحن محتوى بروتينه أقل من دقيق جميع الأغراض. وهو مثالي للمخبوزات الطرية كالكعك والكيك والمعجنات.

\item[الطية الحلزونية \textenglish{(Coil fold)}]
تقنية خاصة من تقنيات الطيّ والشدّ. الطية الحلزونية لطيفة جدًا على العجين، ولهذا فهي ممتازة على امتداد التخمير الأول. فباستخدامها تتحسن قوة العجين مع تقليص الضرر الذي يلحق ببنية العجين إلى أدنى حد.

\item[اللُّبّ \textenglish{(Crumb)}]
القوام الداخلي للخبز، الذي يتسم بحجم الثقوب (الحويصلات) وشكلها وتوزيعها. وهو ما تجده في الداخل حين تفتح رغيف خبز. فاللُّبّ المحكم \textenglish{(\emph{tight crumb})} يعني خبزًا ثقوبه صغيرة موزعة بانتظام، بينما اللُّبّ المفتوح \textenglish{(\emph{open crumb})} له ثقوب أكبر وأكثر عدم انتظام.

\item[المالت الإنزيمي \textenglish{(Diastatic Malt)}]
حبة مالت جُففت ثم طُحنت مسحوقًا. يحتوي هذا المالت إنزيمات تفكك النشا إلى سكريات، وهو ما قد يكون نافعًا في عملية تخمير الخبز. وعند إضافته إلى العجين يمكن أن يحسّن نكهة الخبز ولونه وعمره التخزيني.

\item[بادئة زائدة \textenglish{(Discard)}]
الجزء الذي يُزال من البادئة الحامضة ولا يُغذّى عند صيانة البادئة. يُفعل هذا غالبًا لمنع البادئة من أن تصبح أكبر من اللازم وأصعب في التحكم. ويمكن استخدام البادئة الزائدة في وصفات متعددة أو التخلص منها.

\item[التقسيم \textenglish{(Dividing)}]
عملية تقسيم كتلة العجين إلى قطع أصغر، لتشكيلها عادة إلى أرغفة أو حصص فردية.

\item[ترطيب العجين \textenglish{(Dough Hydration)}]
يُعبَّر عنه كنسبة مئوية، وهو كمية الماء في العجين نسبةً إلى كمية الدقيق. فالعجين ذو الترطيب الأعلى يكون أكثر بللًا ولزوجة، وذو الترطيب الأدنى أكثر صلابة. مثلاً، عجين فيه \LR{\qty{500}{\gram}} دقيق و\LR{\qty{375}{\gram}} ماء ترطيبه \LR{\qty{75}{\percent}}.

\item[قوة العجين \textenglish{(Dough Strength)}]
تشير إلى مرونة العجين وصموده وبنيته. فالعجين القوي يمكن مدّه دون أن يتمزق ويحفظ شكله جيدًا. ويتأثر ذلك بدرجة كبيرة بمحتوى الدقيق من البروتين وبمدى تطوير شبكة الغلوتين.

\item[وعاء خبز هولندي \textenglish{(Dutch Oven)}]
قدر ثقيل بغطاء محكم الإغلاق، مصنوع غالبًا من الحديد الزهر. يُستخدم في الخَبز لاحتجاز البخار خلال الطور الأول من الخَبز، مما يساعد على خلق سطح خارجي مقرمش للخبز.

\item[المرونة \textenglish{(Elasticity)}]
خاصية في العجين تصف قدرته على العودة إلى شكله الأصلي بعد أن يُمَدّ أو يتشوه. تتأثر بمحتوى الدقيق من البروتين وبمدى تطوير شبكة الغلوتين.

\item[القابلية للتمدد \textenglish{(Extensibility)}]
تشير إلى قدرة العجين على أن يُمَدّ أو يُبسَط دون أن يتمزق. وهي عكس المرونة، وهي مرغوبة في أنواع معينة من الخبز، مثل التشباتا ذي بنية اللُّبّ المفتوح.

\item[التغذية \textenglish{(Feed)}]
فعل إضافة دقيق وماء طازجين للحفاظ على البادئة الحامضة. والتغذية المنتظمة تبقي البادئة نشطة ومعافاة.

\item[التخمير \textenglish{(Fermentation)}]
العملية الأيضية التي تحوّل بها الكائنات الدقيقة مثل الخميرة والبكتيريا الكربوهيدرات (مثل السكريات) إلى كحول أو أحماض. وفي صناعة الخبز ينتج عن ذلك ثاني أكسيد الكربون الذي يجعل العجين ينمو.

\item[الفيرمنتولايز \textenglish{(Fermentolyse)}]
استخدام كمية صغيرة من البادئة لإبطاء التخمير. وهي طريقة يجري فيها الجمع بين التخمير والأوتولايز. تُستخدم عادة نحو \LR{\qty{10}{\percent}} من البادئة في الفيرمنتولايز، حيث يُخلط الدقيق والماء والبادئة معًا. فبإضافة البادئة مبكرًا يصير العجين أكثر قابلية للتمدد وأسهل في التعامل.

\item[اختبار النقر بالإصبع \textenglish{(Finger poke test)}]
اختبار النقر بالإصبع طريقة بسيطة وفعالة لمعرفة ما إذا كان خبزك الحامض جاهزًا للخَبز. بعد النمو النهائي، غِرِ إصبعك بالدقيق قليلًا ثم اضغط بلطف نحو نصف بوصة داخل العجين. فإن ارتد العجين ببطء وترك انبعاجًا طفيفًا فقد بلغ الكمال وهو جاهز للفرن. وإن ارتد بسرعة فهو يحتاج مزيدًا من الوقت لينمو. أما إذا انهار العجين أو لم يرتد إطلاقًا فقد يكون مُخمَّرًا بإفراط.

\item[اختبار الطفو \textenglish{(Float test)}]
اختبار الطفو تقنية لتقييم جاهزية البادئة الحامضة. لإجرائه خذ عيّنة صغيرة من بادئتك وضعها بلطف في كوب ماء. ويمكن لنتيجة هذا الاختبار أن تقدم فكرة عن مرحلة تخمير بادئتك.
\begin{itemize}
    \item[] النتيجة الإيجابية: إذا طفت بادئتك بسلاسة على سطح الماء فهذه إشارة واضحة إلى أنها بلغت ذروة تخميرها وجاهزة لتُستخدم عاملًا رافعًا في عجينك. وهذا الطفو نتيجة لغاز ثاني أكسيد الكربون الناتج أثناء عملية التخمير النشط.
    \item[] النتيجة السلبية: بالمقابل، إذا غاصت بادئتك إلى قاع الكوب فهذا يعني أنها ليست جاهزة بعد؛ إشارةً إلى أن عملية التخمير لم تتقدم بما يكفي للحصول على قدرة رافعة مثالية.
\end{itemize}
يستحق الذكر أن اختبار الطفو، مع كونه مؤشرًا موثوقًا للبادئات الحامضية المعتمدة على القمح، قد لا يكون بنفس الفعالية للبادئات غير القمحية. لأن الغاز الناتج عن التخمير في البادئات غير القمحية يميل إلى التسرب بسهولة أكبر، فيقل طفوها. وللبادئات غير القمحية فإن الأسلوب الأدق يتضمن ملاحظة وجود الفقاعات في برطمان بادئتك وتقييم رائحتها؛ فالبادئة الناضجة ينبغي أن تصدر رائحة حامضة رقيقة، لكن دون أن تكون نفّاذة أكثر من اللازم.

\item[لُبّ الأحمق \textenglish{(Fool's Crumb)}]
مصطلح يصف بنية لُبّ تضم عدة جيوب أو ثقوب كبيرة، بدل توزيع متساوٍ للثقوب الأصغر. وليست هذه ميزة مرغوبة بالضرورة، لأنها قد تشير إلى تخمير غير متساوٍ أو تقنيات تشكيل غير سليمة.

\item[الغلوتين \textenglish{(Gluten)}]
معقد بروتيني يتكون من الغليادين والغلوتينين، يوجد في القمح وبعض الحبوب الأخرى. يمنح العجين مرونة وقوة حين يُصفَّى ويُطوَّر كما ينبغي. وعلى مدار التخمير الأول يتدهور كثير من الغلوتين بفعل إنزيم البروتياز وبكتيريا حمض اللبنيك.

\item[التجنيس \textenglish{(Homogenizing)}]
فعل خلق خليط متجانس وموحد. فبالنسبة لدقائق مثل الإيكنورن والجاودار، التي ليس اصطفاف الغلوتين فيها هو الهدف الرئيسي، يضمن العجن أن يبلغ العجين هذه القوامية المتجانسة.

\item[الهُوش \textenglish{(Hooch)}]
طبقة سائلة تتكون أحيانًا فوق البادئة الحامضة. وهي إشارة إلى أن البادئة جائعة وتحتاج تغذية. وتعمل درعًا واقيًا يمنع البادئة من التعرض للعَفَن. ويمكن خلطها من جديد في البادئة أو استخلاصها لصنع الصلصات الحارة.

\item[العجن \textenglish{(Kneading)}]
العملية اليدوية أو الميكانيكية لمواءمة العجين، لتطوير الغلوتين في الخبز المعتمد على القمح والسبلت، أو لتجنيس كتلة العجين في دقائق مثل الإيكنورن أو الجاودار.

\item[القطعة المطبوخة \textenglish{(Kochstück)}]
عند إعداد \textenglish{Kochstück} يُسخَّن الدقيق أو الحبوب مع السائل معًا. ويجب تقليب الخليط أثناء التسخين لمنعه من التكتل أو الاحتراق.

\item[حمض اللبنيك \textenglish{(Lactic Acid)}]
حمض عضوي آخر تنتجه بكتيريا حمض اللبنيك أثناء التخمير. يمنح الخبز الحامض طعمًا حامضًا رقيقًا يشبه الزبادي، ومع حمض الخليك يسهم في حمضية الخبز الإجمالية.

\item[الليفان \textenglish{(Levain)}]
انظر البادئة.

\item[تفاعل مايار \textenglish{(Maillard Reaction)}]
تفاعل مايار أحد أسباب اسمرار الطعام أثناء الطبخ. يحدث التفاعل بين السكريات المُرجِعة والأحماض الأمينية، وحسب المتفاعلات الأولية وظروف الطبخ قد ينتج تنوعًا واسعًا من النواتج النهائية ذات المذاقات والروائح المختلفة. تحدث تفاعلات مايار بسهولة فوق \LR{\qty{150}{\degreeCelsius}}، مع أنها تظل تحدث ببطء أكبر بكثير تحت تلك الحرارة. ويبلغ معدل التفاعل الأمثل بين \LR{pH 6.0} و\LR{pH 8.0}، مع أنه يُفضَّل في الظروف القاعدية.

\item[المالتوز \textenglish{(Maltose)}]
سكر ينتج من التفكك الإنزيمي للنشا بفعل الأميلاز. وهو مصدر غذاء رئيسي للخميرة أثناء التخمير.

\item[المالت غير الإنزيمي \textenglish{(Non-diastatic Malt)}]
حبة مالت جُففت عند درجات حرارة أعلى، فتعطّلت إنزيماتها. يُستخدم أساسًا للنكهة واللون في صناعة الخبز. في العجين الرطب تبقى أميلازات الدقيق نفسها تعمل جيدًا فوق \LR{\qty{50}{\degreeCelsius}}، وتبقى مستقرة حتى نحو \LR{\qty{70}{\degreeCelsius}}، ثم تتعطل مع ارتفاع حرارة العجين أكثر أثناء الخَبز~\cite{wheat+amylase+singh,innograin+enzymes,reichenberger+amylase}.

\item[الانتفاخ في الفرن \textenglish{(Oven Spring)}]
النمو السريع للعجين في الفرن خلال المراحل الأولى من الخَبز، بفعل تمدد الغازات المحتجزة والماء.

\item[الإفراط في التخمير \textenglish{(Over Fermenting)}]
مشكلة شائعة عند عمل عجائن القمح أو السبلت. حين يُخمَّر العجين مدة أطول من اللازم يتفكك معظم الغلوتين فيه. والعجين الناتج لزج جدًا. والخبز النهائي سيكون مسطحًا جدًا ويفقد بعض قوامه المعتاد. وتتميز بنية اللُّبّ فيها بجيوب هوائية دقيقة كثيرة. ويمكن للكثير من الغازات المحتجزة أن يتسرب من العجين أثناء الخَبز. إذا لاحظت هذا أثناء التخمير الأول فمن المستحسن وضع الرغيف داخل قالب خبز ثم خبزه بعد راحة من \LR{30} إلى \LR{60} دقيقة.

\item[الإفراط في التخمير النهائي \textenglish{(Over Proofing)}]
نفس الإفراط في التخمير، لكنه يحدث خلال مرحلة التخمير النهائي.

\item[الرقم الهيدروجيني \textenglish{(pH)}]
مقياس لحموضة المحلول أو قلويتها. يمتد مقياس الرقم الهيدروجيني من 0 إلى 14، حيث القيمة 7 متعادلة. والمحاليل التي يقل رقمها الهيدروجيني عن 7 حمضية، وما فوق 7 قاعدية. والأطعمة المخمرة التي يقل رقمها الهيدروجيني عن 4.2 تُعد آمنة غذائيًا عمومًا. ويمكن استخدام جهاز قياس الرقم الهيدروجيني لمراقبة تقدم تخمير خبزك الحامض.

\item[قيمة \textenglish{P/L} \textenglish{(P/L Value)}]
معيار حاسم مستخلص من اختبار الألفيوغراف؛ تمثل قيمة \textenglish{P/L} النسبة بين ثبات العجين (\textenglish{P}) وقابليته للتمدد (\textenglish{L}). وتحديدًا:
\begin{itemize}
    \item[] \textenglish{P} (الضغط) يشير إلى الضغط اللازم لنفخ العجين أثناء اختبار الألفيوغراف. وهو يدل على مقاومة العجين للتشوه، أي على قوته.
    \item[] \textenglish{L} (الطول) يمثل قابلية العجين للتمدد، أو مدى إمكان مدّه قبل أن يتمزق.
\end{itemize}
وتقدم نسبة \textenglish{P/L} فكرة عن التوازن بين مرونة العجين وقابليته للتمدد:
\begin{itemize}
    \item[] قيمة \textenglish{P/L} منخفضة تعني عجينًا قابلًا للتمدد أكثر من المقاومة. وهذا يعني أن العجين يمكن مدّه بسهولة، مما يجعله مناسبًا لمنتجات معينة مثل البيتزا أو التشباتا.
    \item[] قيمة \textenglish{P/L} مرتفعة تعني عجينًا قوته أكبر من قابليته للتمدد. وهذا العجين أكثر مقاومة للتشوه، وهو ما قد يكون أفضل للمنتجات التي تحتاج حجمًا جيدًا وبنية سليمة، مثل أنواع معينة من الخبز.
\end{itemize}
وتساعد قيمة \textenglish{P/L} الخبازين والمطاحن على تحديد ملاءمة دقيقٍ ما لتطبيقات خبزٍ معينة. ويمكن إجراء تعديلات في خلطات الدقيق أو عمليات الخَبز بناءً على هذه النسبة لتحقيق خصائص الخبز المرغوبة.

\item[البادئة المسبقة \textenglish{(Preferment)}]
خليط من جزءٍ من مكوّنات العجين يُسمح له بالتخمير قبل إضافته إلى عجين الخبز النهائي. ويمكن أن يشمل الخبز الحامض، أو البوليش، أو البيغا، أو \textenglish{pâte fermentée}، أو البادئة الإسفنجية العامة.

\item[التشكيل الأولي \textenglish{(Preshaping)}]
حين تقسم كتلة العجين الكبيرة إلى حصص أصغر تنتهي بقطع عجين غير متجانسة. وهذا يجعل التشكيل أصعب بكثير لأن العجين المشكّل الناتج لن يكون منتظمًا. لهذا السبب يسحب الخبازون قطع العجين الصغيرة فوق سطح منضدة العمل لإنشاء كرات عجين ذات مظهر أكثر انتظامًا.

\item[التخمير النهائي \textenglish{(Proof)}]
النمو الأخير للعجين المشكّل قبل الخَبز.

\item[البروتياز \textenglish{(Protease)}]
إنزيم يفكك البروتينات، بما فيها الغلوتين، إلى سلاسل ببتيدية أصغر وأحماض أمينية. في سياق صناعة الخبز يمكن لنشاط البروتياز أن يكون نافعًا للخبازين ومتحدّيًا لهم معًا. فالنشاط المعتدل للبروتياز يمكن أن يجعل العجين أكثر قابلية للتمدد، وهو ما قد يكون نافعًا في بعض عمليات صناعة الخبز. أما النشاط المفرط فيمكن أن يُضعف شبكة الغلوتين، فينتج عجائن رخوة لزجة يصعب التعامل معها، وقد ينتج عنه خبز بحجم وبنية ضعيفين. وعوامل مثل زمن التخمير ودرجة حرارة العجين ومصدر الدقيق تؤثر في نشاط البروتياز في عجائن الخبز. في الخبز الحامض يمكن لأزمنة التخمير الأطول، خصوصًا عند درجات حرارة أدفأ، أن تؤدي إلى نشاط أعلى للبروتياز، لأن الظروف الحامضة تنشّط بروتيازات الحبوب. ويمكن للدقيق المستخلص من حبوب منبتة أو حبوب مالت أن يحتوي نشاط بروتياز أعلى بفعل عملية النبت أو المالت. إن فهم نشاط البروتياز والتحكم فيه أمر حاسم لبلوغ جودة الخبز المطلوبة وخصائص التعامل المطلوبة.

\item[رغيف بولمان \textenglish{(Pullman Loaf)}]
نوع من أرغفة الخبز يتميز بشكلٍ مستطيل مثالي ولُبّ ناعم فاتح. يُخبز في قالب خاص بغطاء يسمى قالب بولمان أو قالب \textenglish{\emph{pain de mie}}. يضمن الغطاء أن ينمو الخبز بشكل مستقيم تمامًا، بلا القمة القبة المميزة لأرغفة الخبز الأخرى. وكثيرًا ما تُقطَّع أرغفة بولمان شرائح رفيعة جدًا وهي محبوبة لصنع السندويشات.

\item[الإبطاء \textenglish{(Retarding)}]
عملية إبطاء التخمير خلال مرحلة التخمير النهائي بوضع العجين في بيئة أبرد، عادة الثلاجة. يساعد هذا الخبازين على تنظيم الجدول، إذ يمنحهم تحكمًا أكبر في موعد خبز خبزهم، خصوصًا في المخابز واسعة النطاق حيث التوقيت جوهري لخدمة الخبز الطازج لعملاء الصباح الباكر. ومع أن تنظيم الجدول هو السبب الرئيسي، فإن بعض الخبازين يؤكدون أيضًا أن الإبطاء يمكن أن يعزز بروفايل النكهة الإجمالي للخبز. ويُعرف أيضًا بالتخمير في الثلاجة.

\item[الجاودار \textenglish{(Rye)}]
نوع من الحبوب يُستخدم في الخَبز. نظرًا إلى انخفاض محتواه من الغلوتين، فإن الخبز المصنوع من دقيق الجاودار وحده يميل إلى الكثافة. لكن الجاودار ذو نكهة فريدة وفوائد صحية عديدة، لذا فإنه غالبًا ما يُخلط بدقيق القمح في الخَبز. وعادة ما يُصنع خبز الجاودار الخالص بعملية الخبز الحامض لمساعدة العجين على النمو.

\item[السكب المغلي \textenglish{(Scald)}]
طريقة يُسكب فيها الماء المغلي فوق الدقيق أو الحبوب أو مكوّنات أخرى ثم يُترك ليبرد. في الخَبز يمكن لهذه العملية أن تجلتن النشا في الدقيق أو الحبوب، فينتج عجين يحتفظ بالرطوبة أفضل، ويمنح لُبًّا أطرى، وقد يمدّ عمر الخبز التخزيني. إضافة إلى ذلك، يمكن للتسليق أن يساعد على تعطيل إنزيمات معينة قد تضر بجودة العجين. كما يمكن لتقنية التسليق أن تعزز النكهة الإجمالية والرائحة للخبز، فتُبرز نُوتات حبيبية أوضح وتقلل المرارة الموجودة أحيانًا في بعض الحبوب كاملة الحبة.

\item[التشريط \textenglish{(Scoring)}]
قطع سطح عجين الخبز قبل خبزه. يتيح ذلك للعجين أن يتمدد بحرية في الفرن، فلا ينفجر بطرق غير متوقعة. كما يوفر جمالية مضبوطة للرغيف المكتمل.

\item[النخل \textenglish{(Sift)}]
تمرير الدقيق أو مكوّن جاف آخر عبر منخل لإزالة التكتلات وهواءته.

\item[الخلطة المنقوعة \textenglish{(Soaker)}]
خليط من الحبوب أو البذور مع ماء يُترك منقوعًا ليلة كاملة (أو مدة محددة) قبل دمجه في عجين الخبز. يساعد هذا على تلين الحبوب أو البذور (السمسم، بذور القرع، وغيرها) وترطيبها، مما يجعل دمجها في العجين أسهل ويمنح لُبّ الخبز المكتمل رطوبة أكبر.

\item[البادئة الإسفنجية \textenglish{(Sponge)}]
نوع من البادئات المسبقة؛ الإسفنج خليط رطب من الدقيق والماء والخميرة يُترك ليختمر مدة معينة قبل دمجه في العجين النهائي.

\item[البادئة \textenglish{(Starter)}]
خليط مختمر من الدقيق والماء يحتوي مستعمرة من الكائنات الدقيقة تشمل الخميرة البرية وبكتيريا حمض اللبنيك. تُستخدم لرفع الخبز.

\item[العجين المباشر \textenglish{(Straight Dough)}]
طريقة صناعة خبز تُخلط فيها كل المكوّنات معًا دفعة واحدة، دون استخدام بادئة مسبقة.

\item[الطيّات والشدّات \textenglish{(Stretch and Fold)}]
\textenglish{S\&F} تقنية تُستخدم خلال مرحلة التخمير الأول لتقوية العجين والمساعدة في اصطفاف بنية الغلوتين. فبدلاً من العجن التقليدي، يُشَدّ العجين بلطف ثم يُطوى فوق نفسه. وتُكرَّر هذه العملية عادة عدة مرات على مدار التخمير الأول.

\item[التانجونج \textenglish{(Tangzhong)}]
تقنية صينية لصناعة الخبز، تشبه الطريقة اليابانية اليوداني. تتضمن طهي جزء صغير من الدقيق مع الماء (أو الحليب) لإعداد عجينة سائلة أو «رو». وهذه العملية، التي يمكن رؤيتها كصنف من السكب المغلي، تجلتن النشا في الدقيق، فينتج خبز أطرى وأمَلَى مع احتفاظ أفضل بالرطوبة. وبعد أن يبرد، يُخلط التانجونج مع بقية المكوّنات لإنتاج العجين النهائي.

\item[اللُّبّ المحكم \textenglish{(Tight Crumb)}]
يشير إلى لُبّ الخبز (الجزء الداخلي الطري من الخبز) ذي الثقوب الهوائية الصغيرة المتساوية.

\item[الخميرة البرية \textenglish{(Wild Yeast)}]
خميرة موجودة طبيعيًا في البيئة وعلى سطح الحبوب، تُستخدم في تخمير الخبز الحامض بدلاً من الخميرة التجارية. توجد خميرة برية على أي سطح نباتي تقريبًا. وتعيش الخمائر البرية في تكافل مع النبات، فتوفر له درعًا ضد المسببات المرضية وتتلقى مقابل ذلك سكريات من التركيب الضوئي للنبات. وإذا ضعف النبات أصبحت الخمائر البرية طفيلية يمكنها أن تلتهم مضيفها.

\item[قيمة \textenglish{W} \textenglish{(W-Value)}]
معيار يمثل قوة الدقيق من حيث جودة الخَبز. قيمة \textenglish{W}، المستخلصة من اختبار الألفيوغراف شوبان، تقيس الطاقة اللازمة لنفخ فقاعة بالعجين حتى تنفجر. وهي مؤشر مباشر على قدرة الدقيق على تحمل عمليتي التخمير والخَبز. فقيمة \textenglish{W} الأعلى تدل عادة على دقيق أقوى، مناسب للخبوزات عالية الحجم وذات أزمنة تخمير أطول. بالمقابل، تشير قيمة \textenglish{W} الأدنى إلى دقيق أضعف، أنسب للمنتجات التي تحتاج بنية أقل، مثل الكعك والمعجنات.

\item[الخميرة \textenglish{(Yeast)}]
كائنات دقيقة تختمر السكريات الموجودة في العجين، منتجًا الكحول وثاني أكسيد الكربون والحرارة؛ وبذلك يرتفع العجين.

\item[اليوداني \textenglish{(Yudane)}]
طريقة يابانية لصناعة الخبز تتضمن إعداد بادئة بخلط الماء المغلي مع دقيق الخبز بنسبة محددة، عادة \LR{1:1} بالوزن. بعد الخلط يُترك المعجون ليبرد إلى حرارة الغرفة، ثم يُحفظ في الثلاجة ليلة كاملة. وفي اليوم التالي يُخلط مع بقية المكوّنات لصنع العجين. وطريقة اليوداني، وهي في جوهرها صنف من السكب المغلي، تساعد على تحسين قوام الخبز، فتجعله أطرى وأمَلَى، وتعزز أيضًا عمره التخزيني.

\end{description}
